
\documentclass[aspectratio=169]{beamer}
\usetheme{UCLA}
\usepackage{note/lyu}
% ---- Packages you may want ----
\usepackage{amsmath, amsfonts, amssymb}
\usepackage{booktabs}
\usepackage{hyperref}
\usepackage{csquotes}
\usepackage{graphicx}
\usepackage{caption}
\usepackage{listings}
\usepackage[backend=biber,style=authoryear,maxbibnames=99,doi=false,isbn=false,url=true]{biblatex}

\addbibresource{references.bib}
\DeclareCiteCommand{\footcite}[\mkbibfootnote]
  {\usebibmacro{prenote}}
  {\printnames[family-given]{labelname}%
   \setunit{\addnbspace}
   \printfield{journaltitle}%
   \setunit{\addnbspace}
   \printlabeldateextra}
  {\addsemicolon\space}
  {\usebibmacro{postnote}}
\DeclareDelimFormat{multicitedelim}{\addsemicolon\space}
\DeclareMultiCiteCommand{\footcites}[\mkbibfootnote]{\footcite}{\multicitedelim}

% Listings setup
\lstdefinestyle{ucla}{
  basicstyle=\ttfamily\footnotesize,
  keywordstyle=\bfseries\color{uclaBlue},
  commentstyle=\itshape\color{uclaGray},
  stringstyle=\color{uclaDarkBlue},
  frame=single,
  showstringspaces=false,
  breaklines=true,
  tabsize=2
}
\lstset{style=ucla}

% Metadata
\title[Parti2Field]{Parti2Field:Data-Driven Approximation of Mean-Field Limits from Interacting Particle Systems}
% \subtitle{Optional Subtitle}
\author[L. Lyu]{Liyao Lyu}
\cooperator{Hayden Schaeffer, Xinyue Yu}
\institute[UCLA Mathematics]{Department of Mathematics \\ University of California, Los Angeles}
\date{\today}

% Optional logo on title page (replace file with the official UCLA logo if permitted)
\titlegraphic{\hfill 
\includegraphics[height=1.2cm]{ucla_math_logo.png}}

\begin{document}

% Title
\begin{frame}[plain]
  \titlepage
\end{frame}

% Outline
\begin{frame}{Outline}
  \tableofcontents
\end{frame}


% =========================================================
\begin{frame}{Motivation}
\begin{itemize}
  \item We have a \textbf{true} mean-field system (McKean--Vlasov SDE), but its drift is only approximately known.
  \item Model bias and initialization errors cause the approximate system to drift away over time.
  \item We also have \textbf{partial observations} (often of the law / density, not full state).
  \item Goal: add a \textbf{feedback control (nudging)} term so the model distribution stays consistent with observations.
\end{itemize}
\vspace{0.5em}
Key idea: choose nudging as a \textbf{Wasserstein gradient flow} of an observation misfit functional.
\end{frame}

% =========================================================
\section{Problem Setup}

\begin{frame}{True and Model Dynamics (Mean-Field SDE)}
Let $(\Omega,\mathcal{F},\mathbb{P})$ with filtration $(\mathcal{F}_t)_{t\ge 0}$.

\medskip
\textbf{True process} $X_t\in\Rd$:
\[
dX_t = b_{\text{true}}(X_t,\mu_t)\,dt + \Sigma\, dW_t,\qquad X_0\sim \mu_0,
\]
where $\mu_t=\Law(X_t)\in\mathcal{P}_2$.

\medskip
\textbf{Approximate model} $Y_t$:
\[
dY_t = b_{\text{model}}(Y_t,\nu_t)\,dt + \Sigma\, dW_t,\qquad \nu_t=\Law(Y_t).
\]

\medskip
\textbf{Model bias:}\quad
\[
R(x,\mu) := b_{\text{model}}(x,\mu)-b_{\text{true}}(x,\mu).
\]
\end{frame}

% =========================================================
\begin{frame}{Observation Operator}
Assume we observe the law through a smoothed density map:
\[
\mathcal{O}_h(\mu_t)(z) = \int K_h(z-x)\,\mu_t(dx) \;=\; (K_h * \mu_t)(z).
\]

\medskip
In practice: observations may be obtained from partial trajectories/locations
$\{X^{\text{obs},k,M}_t\}_{k=1}^M$ and approximated by empirical measures / KDE.

\medskip
Objective: drive $\nu_t$ toward observation-consistency without knowing $b_{\text{true}}$ exactly.
\end{frame}

% =========================================================
\section{Nudging via Wasserstein Gradient Flow}

\begin{frame}{Quadratic Observation Misfit Functional}
Define the observation misfit cost:
\[
J(\nu,\mu^{\text{obs}}) :=
\frac12 \int_{\Rd} \Big|(K_h*\nu)(x)-(K_h*\mu^{\text{obs}})(x)\Big|^2\,dx.
\]

\medskip
Its first variation (Frech\'et derivative in measure sense):
\[
\frac{\delta J}{\delta \nu}(x)
= \big(\widetilde K_h * (\nu-\mu^{\text{obs}})\big)(x),
\qquad \widetilde K_h := K_h * K_h
\]
(assuming symmetry so convolution is self-adjoint on $L^2$).
\end{frame}

% =========================================================
\begin{frame}{Nudging Vector Field}
In Wasserstein geometry, gradient flow corresponds to velocity field
$-\grad (\delta J/\delta \nu)$.

\medskip
With nudging intensity $\lambda>0$:
\[
u_{\text{nud}}(y,\nu,\mu^{\text{obs}})
:= -\lambda\,\grad_y\!\left(\frac{\delta J}{\delta \nu}(y)\right)
= -\lambda\,\grad\!\big(\widetilde K_h * (\nu-\mu^{\text{obs}})\big)(y).
\]

\medskip
Assimilated mean-field SDE:
\[
dY_t = b_{\text{model}}(Y_t,\nu_t)\,dt
+ u_{\text{nud}}(Y_t,\nu_t,\mu^{\text{obs}}_t)\,dt
+ \Sigma\, dW_t.
\]
\end{frame}

% =========================================================
\section{Particle Approximation}

\begin{frame}{Interacting Particle System}
Approximate $\nu_t$ by particles $\{Y_t^{i,N}\}_{i=1}^N$:
\[
\nu_t^N := \frac1N\sum_{j=1}^N \delta_{Y_t^{j,N}}.
\]

\medskip
Particle dynamics:
\[
dY_t^{i,N} =
b_{\text{model}}(Y_t^{i,N},\nu_t^N)\,dt
-\lambda\Big[
\grad(\widetilde K_h*\nu_t^N)(Y_t^{i,N})
-\grad(\widetilde K_h*\mu_t^{\text{obs}})(Y_t^{i,N})
\Big]dt
+ \Sigma\, dW_t^i.
\]

\medskip
Expanding convolutions (explicit interactions):
\[
dY_t^{i,N} =
b_{\text{model}}(Y_t^{i,N},\nu_t^N)\,dt
-\lambda\left(
\frac1N\sum_{j=1}^N \grad \widetilde K_h(Y_t^{i,N}-Y_t^{j,N})
-\frac1M\sum_{k=1}^M \grad \widetilde K_h(Y_t^{i,N}-X_t^{\text{obs},k,M})
\right)dt
+\Sigma\, dW_t^i.
\]
\end{frame}

% =========================================================
\section{Mean-Field PDE View and Error Dynamics}

\begin{frame}{Nonlinear Fokker--Planck Equation (Formal)}
Let $A=\Sigma\Sigma^\top$.
Writing $r := K_h*(\nu_t-\mu_t^{\text{obs}})$, the assimilated law $\nu_t$ satisfies:
\[
\partial_t \nu_t
= -\grad\cdot(\nu_t\, b_{\text{model}}(x,\nu_t))
+ \grad\cdot(A\grad \nu_t)
+ \lambda\, \grad\cdot\big(\nu_t\,\grad(K_h*r)\big).
\]

\medskip
True mean-field dynamics:
\[
\partial_t \mu_t
= -\grad\cdot(\mu_t\, b_{\text{true}}(x,\mu_t))
+ \grad\cdot(A\grad \mu_t).
\]
\end{frame}

% =========================================================
\begin{frame}{Error Equation}
Define error $v_t := \nu_t-\mu_t$.
Subtracting true from assimilated dynamics yields (schematically):
\[
\partial_t v_t
= -\grad\cdot\Big(\nu_t b_{\text{model}}(x,\nu_t) - \mu_t b_{\text{true}}(x,\mu_t)\Big)
+ \grad\cdot(A\grad v_t)
+ \lambda\,\grad\cdot\big(\nu_t\,\grad(K_h*r)\big),
\]
with $r=K_h*(\nu_t-\mu_t^{\text{obs}})$.

\medskip
In the theorem statement below, a clean case is $\mu_t^{\text{obs}}=\mu_t$ (perfect observations of the law).
\end{frame}

% =========================================================
\section{Stability Result (From the Note)}

\begin{frame}{Assumption Set (Regularity / Stability / Error Sources)}
On periodic domain $\Omega$ (e.g.\ $\mathbb{T}^d$), assume:

\begin{itemize}
  \item[(A1)] Bounded model drift and divergence:
  $\|b_{\text{model}}(\cdot,\rho)\|_{L^\infty}\le B_{\text{model}}$,
  $\|\grad\cdot b_{\text{model}}(\cdot,\rho)\|_{L^\infty}\le B_{\text{model}}$.
  \item[(A2)] Measure-Lipschitz true drift:
  $\|b_{\text{true}}(x,\rho)-b_{\text{true}}(x,\rho')\|\le L_\mu \|\rho-\rho'\|$.
  \item[(A3)] Density bounds: $\|\mu_t\|_{L^\infty}\le \bar\rho$ and $\nu_t(x)\ge \nu>0$ a.e.
  \item[(A4)] Uniform model bias bound:
  $\sup_{t\ge0}\|b_{\text{model}}(x,\rho)-b_{\text{true}}(x,\rho)\|\le \Delta$.
  \item[(A5)] Symmetric kernel $K_h$; convolution self-adjoint and commutes with derivatives.
  \item[(A6)] Smoothing error estimate:
  $\|\grad v - \grad(\widetilde K_h*v)\|_{L^2}\le \delta(h)\|\grad v\|_{L^2}$, with $\delta(h)\to0$ as $h\to0$.
\end{itemize}
\end{frame}

% =========================================================
\begin{frame}{Theorem: Exponential Decay to a Bias Floor}
Let $V(t)=\frac12\|v_t\|^2$ with $v_t=\nu_t-\mu_t$ and assume $\int_\Omega v_0\,dx=0$.
Let $C_P$ be the Poincar\'e constant for zero-mean functions.

\medskip
Assume perfect observation of the law: $\mu_t^{\text{obs}}=\mu_t$.

\medskip
Define the advection-instability coefficient:
\[
C_{\text{adv}} := B_{\text{model}} + \bar\rho\,L_\mu.
\]

Choose $h$ and $\lambda$ so that:
\[
\delta(h)\le \frac{\nu}{2\bar\rho},
\qquad
\lambda > \frac{2}{\nu}\left(\frac{C_{\text{adv}}^2 C_P}{\kappa}-\kappa\right),
\]
where $\kappa$ is the coercivity constant from diffusion.

\medskip
Then there exists $\alpha>0$ such that
\[
\frac{d}{dt}V(t)\le -\alpha V(t) + \frac{\bar\rho^2}{\kappa}\Delta^2,
\qquad
V(t)\le e^{-\alpha t}V(0)+\frac{\bar\rho^2}{\alpha\kappa}\Delta^2.
\]
% Interpretation: exponential contraction until reaching a \textbf{bias floor} $\propto \Delta^2$.
\end{frame}

% % =========================================================
% \begin{frame}{Proof Sketch (Energy Estimate)}
% Start with $V(t)=\frac12\|v_t\|^2$, so $V'(t)=\ip{v_t,\partial_t v_t}$.

% \medskip
% \textbf{Diffusion:}
% \[
% \ip{v_t,\grad\cdot(A\grad v_t)} \le -\kappa\|\grad v_t\|_2^2.
% \]

% \medskip
% \textbf{Advection/model terms:}
% bounded by $C_{\text{adv}}\|\grad v_t\|_2\|v_t\|_2$ plus bias contribution $\bar\rho \|\grad v_t\|_2\Delta$,
% then use Young:
% \[
% C_{\text{adv}}\|\grad v_t\|_2\|v_t\|_2
% \le \frac{\varepsilon}{2}\|\grad v_t\|_2^2 + \frac{C_{\text{adv}}^2}{2\varepsilon}\|v_t\|_2^2.
% \]

% \medskip
% \textbf{Nudging term:} treat as perturbation of identity via smoothing error
% $\|\grad v_t - \grad(\widetilde K_h*v_t)\|\le \delta(h)\|\grad v_t\|$:
% \[
% \text{nudging} \;\lesssim\; -\lambda(\nu-\bar\rho\,\delta(h))\|\grad v_t\|_2^2.
% \]

% \medskip
% Finish with Poincar\'e: $\|\grad v_t\|_2^2 \ge C_P^{-1}\|v_t\|_2^2$ to get $V'(t)\le -\alpha V(t)+(\bar\rho^2/\kappa)\Delta^2$.
% \end{frame}

% =========================================================
\section{Numerical Demonstration}

\begin{frame}{Example: Fish Trajectory Data Assimilation}
Dataset: trajectories of $\approx 1000$ animals (fish) in quasi-2D tank with occlusions / dropouts.

\medskip
Observation density (KDE at discrete times $t_n$):
\[
\mu^{\text{obs}}(x,t_n)=\frac1N\sum_{i=1}^N K_\sigma(x-X_i(t_n)).
\]

\medskip
Pipeline idea:
\begin{itemize}
  \item Train baseline mean-field model drift from data (previous method).
  \item Run the same model with nudging term using observed density.
  \item Compare predicted trajectories/positions vs reference.
\end{itemize}

\medskip
% (Insert your figure here if available.)
\end{frame}

\begin{frame}{Numerical result}
    \begin{figure}
        \centering
        \includegraphics[width=0.5\linewidth]{figure/comparison_2x2_density.png}
        \caption{Caption}
        \label{fig:placeholder}
    \end{figure}
\end{frame}
% =========================================================
\section{Second Model in the Note: Vlasov--Fokker--Planck}

\begin{frame}{From Vlasov to ODEs: the PIC Viewpoint}
\textbf{Vlasov--Poisson} for $f(t,x,v)$ on phase space:
\[
\partial_t f + v\cdot\grad_x f + E(t,x)\cdot\grad_v f = 0,
\qquad
\grad_x\cdot E = 4\pi\rho,\quad \rho(t,x)=\!\int f(t,x,v)\,dv.
\]

\medskip
\textbf{Method of characteristics.} $f$ is constant along curves $(X(t),V(t))$ solving the Hamiltonian ODEs
\[
\frac{dX}{dt}=V,\qquad \frac{dV}{dt}=E(t,X).
\]

\medskip
\textbf{PIC dynamics.} Sample $f_0$ by $N$ macro-particles $\{(X^i_0,V^i_0,w_i)\}_{i=1}^N$ and evolve each by the ODE above:
\[
\dot X^i_t = V^i_t,\qquad
\dot V^i_t = E(t,X^i_t),\qquad
f^N_t \;=\; \sum_{i=1}^N w_i\,\delta_{(X^i_t,V^i_t)}.
\]
Field on a grid is reconstructed from particles: $\rho^N(t,x)=\sum_i w_i\,S(x-X^i_t)$, then solve $\grad_x\!\cdot E^N=4\pi\rho^N$ and interpolate $E^N$ back to particles.

\medskip
\textit{Bridge:} the Vlasov PDE $\Leftrightarrow$ the empirical measure of independent ODE trajectories; adding $\sigma\,dW_t$ to $dV_t$ recovers the SDE on the next slide.
\end{frame}

% =========================================================
\begin{frame}{Vlasov--Fokker--Planck Setting}
Consider phase-space PDE:
\[
\partial_t f + v\cdot \grad_x f + (E(x)+F(t,x))\cdot \grad_v f
= \frac{\sigma^2}{2}\Delta_v f,
\]
with Poisson coupling
\[
\grad_x\cdot E(t,x)=4\pi \rho(t,x),\qquad \rho(t,x)=\int_{\Rd} f(t,x,v)\,dv.
\]

Associated SDE:
\[
dX_t = V_t\,dt,\qquad
dV_t = (E(X_t)+F(t,X_t))\,dt + \sigma\, dW_t.
\]

In this section of the note: model is exact; discrepancy is from imperfect initialization; observation is macroscopic density $\rho(t,x)$.
\end{frame}

% =========================================================
\begin{frame}{Nudged Copy Driven by Density Observation}
Introduce nudged copy:
\[
dY_t = U_t\,dt - u_{\text{nud}}(Y_t,\tilde\rho_t,\rho_t)\,dt,
\qquad
dU_t = (E(Y_t)+F(t,Y_t))\,dt + \sigma\, dW_t,
\]
where $\tilde\rho_t$ is spatial density induced by $Y_t$.

\medskip
Nudging term:
\[
u_{\text{nud}}(t,x)=\lambda\,\grad_x\big(K_h*r(t,\cdot)\big)(x),
\qquad r := K_h*(\tilde\rho-\rho).
\]

Let $\omega := \tilde f - f$ be the phase-space density difference; the note writes the corresponding PDE for $\omega$ including a nudging flux term.
\end{frame}

% =========================================================
\section{Discussion}

\begin{frame}{Takeaways and Open Questions}
\textbf{Takeaways}
\begin{itemize}
  \item Nudging via Wasserstein gradient flow gives an explicit, interpretable feedback control.
  \item Stability: exponential decay in $L^2$ up to a bias floor $\propto \Delta^2$ when $\lambda$ is large enough and $h$ small enough.
  \item Particle implementation is straightforward: add interaction with observed particles / KDE field.
\end{itemize}

\medskip
\textbf{Open questions for group discussion}
\begin{itemize}
  \item How to choose $\lambda$ and $h$ adaptively (bias-variance / stability tradeoff)?
  \item Noisy / partial observations: what changes in the energy estimate?
  \item Finite-$N$ effects: quantitative propagation-of-chaos / concentration for the nudged particle system?
  \item Extensions: other observation operators, constraints, boundaries, and non-symmetric kernels.
\end{itemize}
\end{frame}

% =========================================================
\begin{frame}{References (from the note)}
\small
\begin{thebibliography}{9}
\bibitem{Bessaih2015}
H. Bessaih, E. Olson, and E. S. Titi,
\newblock Continuous data assimilation with stochastically noisy data,
\newblock \emph{Nonlinearity}, 28(3):729, 2015.

\bibitem{WalterCouzin2021}
T. Walter and I. D. Couzin,
\newblock TRex: a fast multi-animal tracking system with markerless identification,
\newblock \emph{eLife}, 10:e64000, 2021.
\end{thebibliography}
\end{frame}



\end{document}
